\section{粘合扩展 Whittaker 范畴，以及谱图景}
\markboth{粘合扩展 Whittaker 范畴，以及谱图景}{}

本节的材料没有经典类比。

我们想描述范畴 $\Whit(G)\ext_{glob}$ 如何由它的子范畴粘合而成，例如 $\Whit(G)\rd_{glob}$ 与 $D(\Bun\rt_T)$，这些子范畴对应 $\Bun^{\mathrm{rat},\mathrm{red},\mathrm{ext}}_N$ 的各个「分层」，而分层由哪些 $\gamma_i\rd$ 取零的模式决定。

粘合数据由这些子范畴之间的函子给出，例如
\[
  (\fI\rt)^!\circ(\fJ\rt)_! : \Whit(G)\rd_{glob}\to D(\overline{\Bun}\rt_T).
\]

\begin{annote}
原文此处把靶写成 $D(\overline{\Bun}\rt_T)$（$\Bun$ 上有横线、同时有 $rat$ 上标）。$(\fI\rt)^!$ 的靶应是 $D(\Bun\rt_T)$（横线表示紧化，与 $rat$ 不能并存）；由 §7.2 的伴随对与 §8.1 的 8.1 式，靶应为 $D(\Bun\rt_T)$。译文照录原文，见工程文件疑点清单。\par\vspace{0.15em}{\raggedleft 【来源：原书 p.16 §8（A）；译者按（类型检查）】\par}
\end{annote}

\ob{8.1}{blk:8.1}\textbf{半单秩 $1$ 的粘合数据。}让我们注意：当 $G$ 的半单秩为 $1$ 时，$\Bun^{\mathrm{rat},\mathrm{red},\mathrm{ext}}_N$ 中「只有两个分层」，即 $\Bun\rt_B$ 与 $\Bun^{\mathrm{rat},\mathrm{red},\mathrm{ext}}_N$，因此范畴 $\Whit(G)\ext_{glob}$ 完全由函子 $(\fI\rt)^!\circ(\fJ\rt)_!$ 描述。

事实上，$D(\Bun^{\mathrm{rat},\mathrm{red},\mathrm{ext}}_N)$ 可以等同于这样的范畴，其对象是三元组 $(\mathcal F_1,\mathcal F_2,\alpha)$，其中
\[
  \mathcal F_1\in\Whit(G)\rd_{glob},\qquad
  \mathcal F_2\in D(\Bun\rt_T),
\]
而
\[
  \alpha\in\Hom_{D(\Bun\rt_T)}\bigl((\fI\rt)^!\circ(\fJ\rt)_!(\mathcal F_1),\mathcal F_2\bigr).
\]
态射按自然的方式取为三元组之间的态射。

对半单秩更高的群，范畴 $\Whit(G)\ext_{glob}$ 由一族这样的函子之间的图来描述，这些函子对应所有互相包含的抛物子群对。

\begin{annote}
「粘合」在这里是「分层—粘合」意义上的：把 $\Bun^{\mathrm{rat},\mathrm{red},\mathrm{ext}}_N$ 按 $\gamma_i\rd$ 取零的模式分成若干层，每层的范畴已知（$\Whit(G)\rd_{glob}$ 对应全非零层，$D(\Bun\rt_T)$ 对应全零层），层与层之间的函子 $(\fI\rt)^!\circ(\fJ\rt)_!$ 就是粘合数据。半单秩 $1$ 时只有两层，故三元组的描述完整；秩高时层数增多，要用「抛物对之间的函子图」。\par\vspace{0.15em}{\raggedleft 【来源：原书 p.16 §8.1（A）；译者按（把分层的含义说清）】\par}
\end{annote}

\ob{8.2}{blk:8.2}\textbf{谱描述。}由推论 6.3.3，我们有
\[
  \Whit(G)\rd_{glob}\simeq
  \Rep(\Ch G,\mathrm{Ran}_X,un)
  \underset{\Rep(\Ch G/[\Ch G,\Ch G],\mathrm{Ran}_X,un)}{\otimes}
  \QCoh\bigl(\LocSys_{\Ch G/[\Ch G,\Ch G]}\bigr),
\]
而由定理 4.6.1，我们有
\[
  D(\Bun\rt_T)\simeq T_{\LocSys_{\Ch B}/\LocSys_{\Ch G}}\text{-}\mathrm{mod}^r .
\]
我们当前的目标是在这些项下把函子 $(\fI\rt)^!\circ(\fJ\rt)_!$ 认出来，即把它认成函子
\begin{equation}
\begin{aligned}
&\Rep(\Ch G,\mathrm{Ran}_X,un)
  \underset{\Rep(\Ch G/[\Ch G,\Ch G],\mathrm{Ran}_X,un)}{\otimes}
  \QCoh\bigl(\LocSys_{\Ch G/[\Ch G,\Ch G]}\bigr)\\
&\qquad\longrightarrow
  T_{\LocSys_{\Ch B}/\LocSys_{\Ch G}}\text{-}\mathrm{mod}^r .
\end{aligned}
\label{eq:8.1}
\end{equation}

下述论断受 V. Drinfeld 提出、A. Arinkin 证明的「Gap 定理」启发。

\begin{othm}{拟定理}{8.2.1}{thm:8.2.1}
粘合函子 \eqref{eq:8.1} 典范同构于复合
\[
\begin{aligned}
&\Rep(\Ch G,\mathrm{Ran}_X,un)
  \underset{\Rep(\Ch G/[\Ch G,\Ch G],\mathrm{Ran}_X,un)}{\otimes}
  \QCoh\bigl(\LocSys_{\Ch G/[\Ch G,\Ch G]}\bigr)\\
&\qquad\xrightarrow{\Loc_{\Ch G,spec}}
  \QCoh(\LocSys_{\Ch G})
  \xrightarrow{(\fP\st)^!}
  T_{\LocSys_{\Ch B}/\LocSys_{\Ch G}}\text{-}\mathrm{mod}^r .
\end{aligned}
\]
\end{othm}

在 [GL:funct] 中我们将说明，定理 8.2.1 如何允许我们对 $G=GL_2$ 构造一个单方向的函子
\[
  \Psi_G : \QCoh_N(\LocSys_{\Ch G})\to D(\Bun_G).
\]

\begin{annote}
Gap 定理（间隙定理）是 Drinfeld 提出、Arinkin 证明的关于 $\Bun_G$ 上 D-模与谱侧之间「间隙」的定理；原文只说本文的论断受其启发，并未用其内容。这里的关键是：粘合函子 $(\fI\rt)^!\circ(\fJ\rt)_!$ 本来是用几何方式定义的，而 \eqref{eq:8.1} 把它翻译到谱侧，认成「先 $\Loc_{\Ch G,spec}$、再 $(\fP\st)^!$」。这一步是 §8 的全部内容，也是 [GL:funct] 构造 $\Psi_G$ 的出发点。\par\vspace{0.15em}{\raggedleft 【来源：原书 p.17 §8.2（A）；译者按（把该论断在全文中的位置说清）】\par}
\end{annote}
